% =====================================================================
%  Définitions du glossaire et des acronymes
%  (référencés dans le texte avec \gls{clé} ou \acrshort{clé})
% =====================================================================

% ---------- Acronymes ------------------------------------------------
\newacronym{ci}{CI}{Continuous Integration (intégration continue)}
\newacronym{cd}{CD}{Continuous Delivery / Deployment}
\newacronym{sdk}{SDK}{Software Development Kit}
\newacronym{api}{API}{Application Programming Interface}
\newacronym{abi}{ABI}{Application Binary Interface}
\newacronym{ide}{IDE}{Integrated Development Environment}
\newacronym{tdd}{TDD}{Test-Driven Development}
\newacronym{poc}{PoC}{Proof of Concept}
\newacronym{roi}{ROI}{Return On Investment (retour sur investissement)}
\newacronym{qa}{QA}{Quality Assurance (assurance qualité)}
\newacronym{vcs}{VCS}{Version Control System}
\newacronym{e2e}{E2E}{End-to-End (test de bout en bout)}

% ---------- Termes ---------------------------------------------------
\newglossaryentry{dette-technique}{
  name={dette technique},
  description={Coût implicite des retouches futures induit par le choix
  d'une solution rapide plutôt que d'une approche meilleure mais plus
  longue à mettre en œuvre. Se rembourse par le refactoring ; génère des
  « intérêts » sous forme de ralentissement du développement}
}
\newglossaryentry{test-regression}{
  name={test de non-régression},
  description={Test vérifiant qu'une modification du code n'a pas
  réintroduit de défaut ni cassé une fonctionnalité existante}
}
\newglossaryentry{test-integration}{
  name={test d'intégration},
  description={Test validant le bon fonctionnement conjoint de plusieurs
  composants assemblés, par opposition au test unitaire}
}
\newglossaryentry{fork}{
  name={fork},
  description={Divergence d'une base de code en plusieurs variantes
  maintenues en parallèle, partageant une origine commune}
}
\newglossaryentry{retro-ingenierie}{
  name={rétro-ingénierie},
  description={Analyse d'un système ou binaire existant pour en
  comprendre la structure et le comportement, en l'absence de
  documentation ou de sources complètes}
}
\newglossaryentry{toolchain}{
  name={chaîne de compilation (toolchain)},
  description={Ensemble des outils (compilateur, éditeur de liens,
  bibliothèques) permettant de transformer le code source en binaire
  exécutable pour une cible donnée}
}
